\paragraph{A detector identifies the model's propagating Weyl flavors.}
The massless sector now has a concrete model-particle readout.  On the
periodic $8^3$ lattice at $c=1/8$, retain all 1,024 physical CAR modes of
the original Weyl/contact action.  For each $b\in\{0,1\}^3$, choose
momentum labels $\ell_b=(4b_x,4b_y,4b_z+1)$, with $k_b=\pi\ell_b/4$.
For component eigenvalue $r=\pm1$ of $\sigma_z$, the normalized mode is
\[
 \psi_{b,r}(x,s)=\frac{e^{ik_b\cdot x}}{\sqrt{512}}
                 \delta_{s,(1-r)/2},\qquad
 H\psi_{b,r}=E_{b,r}\psi_{b,r},\qquad
 E_{b,r}=-r(-1)^{b_z}\frac{\sqrt2}{16}.
\]
The contact term is exactly zero on the entire one-particle sector for
every supplied $U\ge0$.  These sixteen modes are orthonormal, checked
in Gaussian integers with amplitudes $(A+\sqrt2 B)/32$.
Their rank-one projectors satisfy $H\Pi_{b,r}=E_{b,r}\Pi_{b,r}$;
the remaining one-particle projector commutes with $H$.
Thus a label gives an actual source-generator energy, not a fitted
interpretation of a peak.  All modes have $N=1$, odd CAR parity and
fermionic exchange statistics from the supplied CAR algebra.

The operational species are eight Weyl flavors indexed by measured
valley momentum.  Their common energy magnitude is
$\sqrt2/16\simeq0.0883883476483$ in original model units; it alone does
not distinguish the flavors.  With
$\chi_b=-(-1)^{b_x+b_y+b_z}$, the positive-energy branch has cone helicity
$\chi_b/2$ and the negative branch has the opposite helicity.
\begin{center}
\begin{tabular}{c c c c}
flavor $b$ & momentum labels $\ell_b$ & positive-branch component
          & positive cone helicity\\\hline
000 & $(0,0,1)$ & 1 & $-1/2$\\
001 & $(0,0,5)$ & 0 & $+1/2$\\
010 & $(0,4,1)$ & 1 & $+1/2$\\
011 & $(0,4,5)$ & 0 & $-1/2$\\
100 & $(4,0,1)$ & 1 & $+1/2$\\
101 & $(4,0,5)$ & 0 & $-1/2$\\
110 & $(4,4,1)$ & 1 & $-1/2$\\
111 & $(4,4,5)$ & 0 & $+1/2$
\end{tabular}
\end{center}
Here helicity uses the explicit proper valley frame
$R_b=-\chi_b\operatorname{diag}((-1)^{b_j})$; it is
$(-1)^{b_x+b_y}r/2$ on this family.  The raw detector measures momentum
and CAR component, so bare $\sigma_z/2$ is not silently identified with
this cone spin or with physical particle spin.  The total directional
current $J_z(k)=-c\cos k_z\,\sigma_z$ has eigenvalue $E_{b,r}$ here,
with nonzero current $\pm\sqrt2/8192$ on every forward z bond.
Current is stationary on these selected modes; $[H,J_z]\ne0$ at generic
momenta.  The witness is propagating flux, even for a uniform-density
plane wave.

\paragraph{The Fourier readout is compiled with complete native errors.}
The detector applies the fermionic extension of
$F[k,x]=e^{-2\pi i k\cdot x/8}/\sqrt{512}$ separately to each component,
then copies all 1,024 occupations to fresh pointers.  Its output includes
every binary occupation string, the vacuum, all other one-particle
addresses, multiple occupations and every export abort.  On the chosen
sixteen effects it measures the energies above; the other 1,008
one-particle records remain present without that energy assignment.
On higher-number inputs these same records measure mode occupations,
not the total interacting energy.

An explicit three-axis FFT uses 4,608 two-mode butterflies and 768
bit-reversal transpositions.  Each butterfly is a mode phase followed by
the native hopping word with coefficient $-i$ and duration $\pi/4$.
All nonadjacent operations use graded FSWAP routes, including their
double-occupation sign.  Exact arithmetic in
$\mathbb Z[\zeta_8]/(\zeta_8^4+1)$ proves the 64 entries and 64 Gram
sums of the one-axis transform; tensoring gives the full mode map.
One transform requests 1,260,800 adjacent FSWAPs and 32,909,824 native
pulses.  Native excitation from the vacuum, inverse-transform
preparation, forward-transform readout and all occupation pointers total
65,822,721 formal pulses.  Original-$H$ propagation is a separate
specified model-law service, not an uncounted synthesized field word.

The complete literal-source error for preparation and detection is
\[
 \epsilon_{\rm source}<8.31\times10^{-9}.
\]
The estimate retains all first-export failures and is valid with untouched
references.  It does not assume that literal native maps are exactly
Gaussian or number preserving.  Unitary original-$H$ propagation on the
full Fock space does not amplify preparation error.  Therefore, at every
model time, a prepared catalog member produces its correct native label
with unconditional probability at least $1-\epsilon_{\rm source}$;
all wrong, outside, vacuum, multiple-occupation and abort outputs together
have probability at most $\epsilon_{\rm source}$.

The replay executes the ideal Gaussian comparison on all 1,024
one-particle rows and all sixteen inputs.  At time eight its complete
output-vector discrepancy from the exact spectral prediction is about
$1.94\times10^{-15}$, separate from the rigorous native bound.
A seeded 65,536-record illustration retains all sixteen label counts
and every zero outside count for the exact uniform input mixture.
These are model-simulated counts, not a literal $2^{1024}$-row native
execution or laboratory data.  No survival normalization is used.

\paragraph{Mass readout and comparison with measured particles.}
For the periodic family $L=8\cdot2^m$, $a=2^{-m}$ and
$c=(1/8)/a$, the model period stays eight and the momentum offset stays
$p_z=\pi/4$.  The exact branch energies satisfy
\[
 E_\pm(a)=\pm\frac{\sin(a\pi/4)}{8a}
       \longrightarrow\pm\frac{\pi}{32},
 \qquad m_{\rm continuum}=0.
\]
This is a direct finite-torus consequence of the original spectrum.
The finite-spacing mass-shell defect is a regulator correction, not
a tachyonic mass.  No SI scale is assigned.

A separate retrospective comparison uses fixed PDG 2025 targets, which
never enter the source or detector calculation:
\begin{center}
\begin{tabular}{c r c}
physical target & measured mass in $\mathrm{MeV}/c^2$
                & present continuum mass\\\hline
electron & $0.51099895000\pm0.00000000015$ & 0\\
muon & $105.6583755\pm0.0000023$ & 0\\
tau & $1776.93\pm0.09$ & 0
\end{tabular}
\end{center}
No common positive unit conversion maps zero to these positive masses;
equal finite-momentum energies are not a rest-mass hierarchy.\footnote{PDG
2025 listings:
\url{https://pdg.lbl.gov/2025/listings/rpp2025-list-electron.pdf};
\url{https://pdg.lbl.gov/2025/listings/rpp2025-list-muon.pdf};
\url{https://pdg.lbl.gov/2025/listings/rpp2025-list-tau.pdf}.}
The photon has spin one, so direct spin-preserving identification with
these cone-helicity-$1/2$ modes also fails.  Three such massless flavors
give zero mass-squared differences, outside the nonzero ranges of the
specified three-neutrino fit in PDG Table 14.7.\footnote{The photon
listing is \url{https://pdg.lbl.gov/2025/listings/rpp2025-list-photon.pdf}.
The selected normal-ordering fit gives
$\Delta m_{21}^2\in[6.92,8.05]\times10^{-5}\,\mathrm{eV}^2/c^4$ and
$\Delta m_{32}^2\in[2.388,2.531]\times10^{-3}\,\mathrm{eV}^2/c^4$;
these are reported $3\sigma$ fit ranges, not exact bounds:
\url{https://pdg.lbl.gov/2025/reviews/rpp2025-rev-neutrino-mixing.pdf}.}
A single massless neutrino or a composite photon in another sector is
not excluded.  The model's Weyl flavors are now operationally identified;
direct identifications with measured physical species remain zero.
The empty reference is not a ground state, and the negative branch is
retained rather than renamed as a physical antiparticle.  Interacting
vacuum stability, mass generation, physical charges and the empirical
dictionary remain additional work.
